% !TEX root = ../notes.tex

\documentclass[../notes.tex]{subfiles}

\begin{document}

\section{April 29}
Here we go.

\subsection{The Exponents}
Let's see what else we can say about $\mathrm H^\bullet(G;\CC)$, where $G$ is a compact simply connected almost simple Lie group.
\begin{notation}
	Fix a compact simply connected almost simple Lie group $G$. Then let $(m_1,\ldots,m_r)$ be the increasing sequence of positive integers for which the generators of $\mathrm H^\bullet(G;\CC)$ as an exterior Hopf algebra have degree $2m_i+1$.
\end{notation}
\begin{remark} \label{rem:number-of-gens}
	In \Cref{prop:number-of-gens-cohomology}, we showed that $r$ is the rank.
\end{remark}
\begin{example}
	On the homework, we will show that $\left(\land^3\mf g^*\right)^\mf g$ is one-dimensional. Thus, $m_1=1$, and $m_i>1$ for all $i\ge2$.
\end{example}
\begin{remark}
	Note that
	\[\sum_{i=1}^r(2m_i+1)=\dim\mf g,\]
	so $\sum_im_i=\left|\Phi_+\right|$.
\end{remark}
For our next calculations, it will be helpful to have more notation.
\begin{definition}[Poincar\'e polynomial]
	Fix a smooth manifold $M$. Then we define the \textit{Poincar\'e polynomial} $P_M$ to be
	\[P_M(q)\coloneqq\sum_{j\ge0}\dim\mathrm H^j(M;\CC)q^j.\]
\end{definition}
\begin{example}
	One can calculate that $P{\op{SU}(2)}(q)=1+q^3$ and $P_{\op{SU}(3)}(q)=\left(1+q^3\right)\left(1+q^5\right)$ and $P_{\op{Sp}_4}(q)=\left(1+q^3\right)\left(1+q^7\right)$ and $P_{G_2}(q)=\left(1+q^3\right)\left(1+q^{11}\right)$.
\end{example}
\begin{remark}
	By Poincar\'e duality, one can find that $m_{r+1-i}=h-m_i$, where $h$ is the Coxeter number. Indeed, the point is that $\dim G=h(r+1)$. This can be used to compute the $m_i$s in rank $3$.
\end{remark}
In fact, one has the following.
\begin{theorem} \label{thm:get-exps}
	Fix a compact simply connected almost simple Lie group $G$. Then the numbers $\{m_i\}$ are the exponents of the Lie algebra $\mf g$.
\end{theorem}
We will sketch arguments for classical types. For now, we will content ourselves with arguments in type $A$, and we will suggest how to argue in the other classical types, leaving details to the homework.
\begin{proposition} \label{prop:poincare-un}
	We have
	\[P_{\op U(n)}(q)=(1+q)\left(1+q^3\right)\cdots\left(1+q^{2n-1}\right).\]
\end{proposition}
\begin{proof}
	We use Schur--Weyl duality. Let $V$ be the standard representation. Then $\mf g^*=V\otimes V^*$, so by \Cref{cor:skew-howe}, we find that
	\[\land^\bullet(V\otimes V^*)=\bigoplus_{\lambda=(\lambda_1,\ldots,\lambda_n)}\left(S^\lambda V\otimes S^{\lambda^\dagger}V^*\right).\]
	Notably, this has no $\mf g$-invariants if $\lambda\ne\lambda^\dagger$ because $S^\lambda V$ is just some irreducible representation. However, even if $\lambda=\lambda^\dagger$, then we only have a one-dimensional space of invariants by Schur's lemma, so it follows that
	\[P_{\op{U}(n)}(q)=\sum_{\substack{\lambda=(\lambda_1,\ldots,\lambda_n)\\\lambda=\lambda^\dagger}}q^{\left|\lambda\right|}.\]
	Thus, we see that we are hunting for Young diagrams $\lambda$ which are equal to their transpose. Well, the largest such diagram is the full square of length $n$, and we see that the desired Young diagrams are simply obtained by removing one of the ``hooks'' which are a top-left segment. Looping over all decisions to either keep or remove a hook, we find that
	\[P_{\op{U}(n)}(q)=(1+q)\left(1+q^3\right)\cdots\left(1+q^{2n-1}\right),\]
	as required.
\end{proof}
We are now ready to prove \Cref{thm:get-exps} in type $A$, which amounts to the following.
\begin{corollary}
	We have
	\[P_{\op{SU}(n)}(q)=\left(1+q^3\right)\left(1+q^5\right)\cdots\left(1+q^{2n-1}\right).\]
\end{corollary}
\begin{proof}
	Note that $\op U(n)$ contains the subgroup $\op{SU}(n)$, and it is the kernel of the determinant map. Thus, there is a fiber sequence
	\[\op{SU}(n)\to\op U(n)\to\op U(1).\]
	Taking Poincar\'e polynomials reveals that $P_{\op U(n)}(q)=P_{\op U(1)}(q)P_{\op{SU}(n)}(q)$, so it follows that
	\[P_{\op{SU}(n)}=(1+q)^{-1}P_{\op U(n)}(q),\]
	so we are done by \Cref{prop:poincare-un}.
\end{proof}
Here is an application of \Cref{thm:get-exps}.
\begin{theorem}[Macdonald's identity]
	Fix a compact simply connected almost simple Lie group $G$ with Lie algebra $\mf g$. Then we have
	\[\frac{(1+q)^r}{\left|W\right|}\int_T\prod_{\alpha\in\Phi}(1+qe^\alpha)(1-e^\alpha)=\prod_{i=1}^r\frac{1+q^{2m_i+1}}{1+q}.\]
\end{theorem}
\begin{proof}
	The root decomposition implies that $\land^\bullet\mf g^*=\land^\bullet\mf h^*\otimes\bigotimes_\alpha\land^\bullet\mf g_\alpha^*$. Now, the character from $\land^\bullet\mf h^*$ is $(1+q)^r$, and the character from $\land^\bullet\mf g_\alpha^*$ is $(1+qe^\alpha)$, so we find that our Poincar\'e polynomial is
	\[\op{ch}\land^\bullet\mf g^*\cdot q^j=(1+q)^r\prod_{\alpha\in\Phi}\left(1+qe^\alpha\right).\]
	Thus, arguing exactly as in \Cref{prop:number-of-gens-cohomology}, our Poincar\'e polynomial is
	\[\frac{(1+q)^r}{\left|W\right|}\int_T\prod_{\alpha\in\Phi}(1+qe^\alpha)(1-e^\alpha),\]
	so \Cref{thm:get-exps} is equivalent to
	\[\frac{(1+q)^r}{\left|W\right|}\int_T\prod_{\alpha\in\Phi}(1+qe^\alpha)(1-e^\alpha)=\prod_{i=1}^r\frac{1+q^{2m_i+1}}{1+q},\]
	as desired.
\end{proof}
\begin{remark}
	It is also possible to prove Macdonald's identity directly via some calculations, which then proves \Cref{thm:get-exps} by the above argument.
\end{remark}
While we're here, we give another approach to \Cref{thm:get-exps}, which will work for the other classical types. We use cell decompositions.
\begin{proof}[Proof of \Cref{prop:poincare-un}]
	Recall that $\op U(n)$ acts transitively on $S^{2n-1}$ with stabilizer $\op U(n-1)$. On the other hand, $S^{2n-1}$ admits a cell decomposition into cells of dimension $0$ and a cell of dimension $2n-1$, which can be described as viewing $S^{2n-1}$ as the one-point compactification of $S^{2n-1}$. Thus, we can use the fibration to inductively give $\op U(n)$ a cell decomposition into $2^n$ cells.

	Thus, we could compute the cohomology of $\op U(n)$ using cellular cohomology: indeed, the above cells give a cell decomposition. Observe that if the relevant chain complex had any interesting differentials, then $\dim\mathrm H^*(\op U(n);\CC)$ would be smaller than $2^n$, which contradicts \Cref{prop:number-of-gens-cohomology}. Thus, we must have an isomorphism of graded abelian groups
	\[\mathrm H^*(\op U(n);\ZZ)\cong\land^\bullet_\ZZ(\xi_1,\ldots,\xi_n),\]
	where $\deg\xi_i=2i-1$. This must also be an isomorphism of rings for degree reasons; namely, we see that we must have a generator in degree one, but then we cannot access degree three, so there must be a generator in degree three, and we can continue this process inductively.
\end{proof}
\begin{remark}
	On the homework, we will repeat this argument for the symplectic groups. Namely,
	\[\mathrm H^*(\op U_n(\HH);\ZZ)\cong\land^\bullet(\xi_1,\ldots,\xi_n),\]
	where $\deg\xi_i=4i-1$. The moral is that $\op U_n(\HH)/\op U_{n-1}(\HH)\cong S^{4n-1}$.
\end{remark}
\begin{remark}
	There is no analogous proof for the orthogonal groups. Indeed, we will need more cells: these cells can be seen in the cohomology of $\ZZ/2\ZZ$. Instead, one can use the cohomology of Stieffel manifolds, which will be done on the homework.
\end{remark}

\subsection{The Cohomology of Homogeneous Spaces}
Fix a compact connected Lie group $G$, and let $K\subseteq G$ be some closed Lie subgroup; let $\mf g$ and $\mf k$ be the corresponding Lie algebras. We are now interested in computing the cohomology of $G/K$. Because $G$ acts on $K$, its cohomology is computed by the complex $\Omega^\bullet(G/K)^G$ by \Cref{rem:pass-to-group-invariants}. On the other hand, a $G$-invariant differential form on $G/K$ can be spread out from the identity uniquely, so it is a differential form on $\mf g/\mf k$. However, to extend to $G/K$, it must further be invariant by the stabilizer $K$ at the identity, so we conclude that there is an isomorphism
\[\left(\land^\bullet(\mf g/\mf k)^*\right)^K\cong\Omega^\bullet(G/K)^G\]
of complexes. Note that taking $K$-invariants may be done in two steps: we can take $\mf k$-invariants (which is algebraic) and then take invariants for the finite group $K/K^\circ$.

We are now interested in the differential of our complex.
\begin{proposition}
	Fix an inclusion $\mf k\subseteq\mf g$ of Lie algebras. Then there is an inclusion
	\[\left(\land^\bullet(\mf g/\mf k)^*\right)^{\mf k}\subseteq\land^\bullet\mf g^*\]
	of complexes.
\end{proposition}
\begin{proof}
	This will be shown on the homework.
\end{proof}
This motivates the following definition.
\begin{definition}[relative Chevalley--Eilenberg complex]
	Fix an inclusion $\mf k\subseteq\mf g$ of Lie algebras. Then the \textit{relative Chevalley--Eilenberg complex} is the subcomplex $\left(\land^\bullet(\mf g/\mf k)^*\right)^{\mf k}$ of $\land^\bullet\mf g^*$. Its cohomology is denoted $\mathrm H^*(\mf g;\mf k)$. We denote the relative complex by $\op{CE}^*(\mf g;\mf k)$.
\end{definition}
\begin{remark}
	It turns out that the complex $\Omega^\bullet(G/K)$ consists of the $K/K^\circ$-invariants of the complex $\op{CE}(\mf g;\mf k)$. In particular, the spaces are the same, and the differentials agree, which can be checked via \Cref{prop:cartan-diff-formula}. Thus,
	\[\mathrm H^\bullet(G/K;\CC)=\mathrm H^\bullet(\mf g;\mf k)^{K/K^\circ}.\]
	Here, note that $K$ normalizes $K^\circ$, so $K/K^\circ$ does in fact act on $G/K^\circ$.
\end{remark}
Here is a useful trick, which explains why we insist on allowing $K$ to be disconnected.
\begin{lemma} \label{lem:k-has-minus-one}
	Fix a compact connected Lie group $G$, and let $K\subseteq G$ be a closed Lie subgroups; let the Lie algebras be $\mf k\subseteq\mf g$. Suppose there is $\varepsilon\in K$ which acts by $-1$ on $\mf g/\mf k$. Then the cohomology of $G/K$ is supported in even degrees, and
	\[\mathrm H^\bullet(G/K;\CC)\cong\mathrm H^\bullet(\mf g;\mf k)^{K/K^\circ}.\]
\end{lemma}
\begin{proof}
	We may pass to the relative Chevalley--Eilenberg complex. Having an automorphism $\varepsilon$ acting by $-1$ implies that the complex $\left(\land^\bullet(\mf g/\mf k)^*\right)^K$ is supported in the even degrees. Thus, the differential vanishes automatically, and the result follows.
\end{proof}
% \begin{remark}
% 	Here is a useful trick: if $\varepsilon\in K$ acts by $-1$ on $\mf g/\mf k$, then $\left(\land^\bullet(\mf g/\mf k)^*\right)^K$ is supported in the even degrees. Thus, the differential vanishes automatically, and we find that
% 	\[\mathrm H^\bullet(G/K;\CC)\cong\mathrm H^\bullet(\mf g;\mf k)^{K/K^\circ}\cong\left(\land^\bullet(\mf g/\mf k)^*\right)^K\]
% 	already!
% \end{remark}
\begin{exe}
	Fix positive integers $m$ and $n$. We compute the Poincar\'e polynomial of the complex Grassmannian $\op{Gr}_{m+n,n}(\CC)$.
\end{exe}
\begin{proof}
	Fix Hermitian spaces $V$ and $W$ of dimensions $m$ and $n$, and define $G\coloneqq\op U(V\oplus W)$ and $K\coloneqq\op U(V)\times\op U(W)$. Then $G/K=\op{Gr}_{m+n,n}(\CC)$, which classifies $n$-dimensional Hermitian subspaces of $\CC^{m+n}$. Note that $\varepsilon=\begin{bsmallmatrix}
		1 \\ & -1
	\end{bsmallmatrix}$ acts by $-1$ on $\mf g/\mf k$: indeed, $\mf g/\mf k$ is given by matrices of the form $\begin{bsmallmatrix}
		& B \\ C
	\end{bsmallmatrix}$. Thus, by \Cref{lem:k-has-minus-one},
	\[\mathrm H^\bullet(\op{Gr}_{m+n,n}(\CC);\CC)=\land^\bullet\left(V\otimes W^*\oplus W\otimes V^*\right)^{K}.\]
	By \Cref{cor:skew-howe}, we find that this is
	\[\bigoplus_{\lambda,\mu}S^\mu V\otimes S^{\mu^\dagger}W^*\otimes S^\lambda W\otimes S^{\lambda^\dagger}V^*,\]
	where $\lambda$ has at most $n$ parts, and $\mu$ has at most $m$ parts.
	
	Now, taking $K$-invariants, Schur's lemma tells us that this is $\bigoplus_\lambda\CC\left[2\left|\lambda\right|\right]$, where $\lambda$ has at most $n$ parts, and $\lambda^\dagger$ has at most $m$ parts. Accordingly, let $N_i(n,m)$ be the number of Young diagrams with $i$ boxes and which fit into the rectangle of height $n$ and length $m$, so we see that $P_{\op{Gr}_{m+n,n}(\CC)}(q)$ is the generating function $\sum_iN_i(n,m)q^{2i}$. To this end, we let $p_i$ denote the jump $\lambda_i-\lambda_{i+1}$, where $p_0=n-\lambda_1$ and $p_n=\lambda_n$. Then the $p$s are all nonnegative can be seen to uniquely recover $\lambda$, so we may compute with these variables instead: note $\left|\lambda\right|=0p_0+p_1+2p_2+\cdots+np_n$, so
	\[\sum_mf_{n,m}(q)z^m=\sum_{p_0,\ldots,p_m\ge0}q^{2(0p_0+1p_1+2p_2+\cdots+mp_m)}z^{p_0+\cdots+p_m}=\prod_{j=0}^n\frac1{1-q^{2j}z},\]
	which we call good enough for our calculation.
\end{proof}
\begin{example}
	We find that
	\[\sum_mf_{1,m}(q)z^m=\frac1{(1-z)\left(1-q^2z\right)}.\]
	This recovers the Poincar\'e polynomial of $\CP^m$, once we isolate the $z^m$ term.
\end{example}
\begin{remark}
	In general, one has
	\[f_{m,n}(q)=\frac{[m+n]_{q^2}!}{[m]_{q^2}![n]_{q^2}!},\]
	where $[n]_{q^2}=\frac{1-q^{2n}}{1-q^2}$ and $[n]_{q^2}!=[1]_{q^2}\cdots[n]_{q^2}$. Indeed, this follows from the $q$-binomial theorem, which is
	\[\sum_{m\ge0}\frac{[m+n]_{q^2}!}{[m]_{q^2}![n]_{q^2}!}z^m=\prod_{j=0}^n\frac1{1-q^{2j}z}.\]
	Indeed, call the right-hand function $F(z)$, which satisfies
	\[F\left(q^2z\right)=\frac{1-z}{1-q^{2(n+1)}z}\cdot F(z).\]
	One can check that the left-hand side also satisfies this, which completes the proof because there is only one power series in $z$ with this property.
\end{remark}
It is also possible to compute a cell decomposition of the Grassmannian via Schubert cells. Here is the idea, akin to \Cref{lem:k-has-minus-one}.
\begin{remark} \label{rem:even-cells}
	Suppose that a connected compact manifold $X$ admits a cell decomposition with only even-dimen\-sional cells; for example, such things frequently arise from complex manifolds in algebraic geometry. This means that the differential in the cellular chain complex vanishes for degree reasons, so $\mathrm H^\bullet(X;\ZZ)$ admits a basis given by those cells.
\end{remark}
\begin{exe}
	Fix positive integers $m$ and $n$. We compute the integral cohomology of the complex Grassmannian $\op{Gr}_{m+n,m}(\CC)$.
\end{exe}
\begin{proof}
	Choose a complete flag $\{F_i\}$ of $\CC^{m+n}$, such as the standard flag. Now, for any $V\in\op{Gr}_{m+n,m}(\CC)$, the intersections $\{V\cap F_i\}_i$ have dimensions which increase by at most one and rise from $0$ to $m$. Thus, for each $i$, we may find the smallest integer $\ell_i$ for which $\dim(V\cap F_{\ell_i})=i$; for example, $\ell_0=0$. Note that $\{\ell_i\}_i$ is a strictly increasing sequence. Then we may define $\lambda_1\coloneqq\ell_m-m$ and $\lambda_2\coloneqq\ell_{m-1}-(m-1)$ and so on to $\ell_m\coloneqq\ell_1-1$, and we see that $\lambda$ is a partition with $m$ parts each of size at most $n$, meaning that $\lambda$ fits into the rectangle with $m$ rows and $n$ columns.

	Now, for such a partition $\lambda$ fitting into this rectangle, we define $S_\lambda\subseteq\op{Gr}_{m+n,m}(\CC)$ to be the subset of subspaces $V$ which induce the partition $\lambda$. On the homework, we will show that $S_\lambda$ is a locally closed embedded complex submanifold of complex dimension $\left|\lambda\right|$. Thus, $\op{Gr}_{m+n,m}(\CC)$ is covered by even cells, so we find that $\mathrm H^{2i}(\op{Gr}_{m+n,m}(\CC);\ZZ)$ is free over $\ZZ$ with basis given by the Schubert cells.
\end{proof}
\begin{remark}
	For example, it follows that the Grassmannian is simply connected.
\end{remark}
\begin{remark}
	There is a story for Grassmannians over $\RR$, but it is more complicated.
\end{remark}

\end{document}